\def\ChapterCode{ADVANCED}
\chapter
[\CO{[ADVANCED]} Artikelsammlung]
{Artikelsammlung}
\label{chp:ADVANCED-Artikel}

\ChapterIntro
{%

}
{%
\begin{description}
  \item[Maintainer:] \ldots
  \item[Autoren:] Diverse
  \item[Reviewer:] Standard-Reviewprozess
  \item[Version:] {\VarDocVersion} ({\VarDocVersionDescription})
  \item[SHA1:] {\DocGitCommit}
\end{description}

{\TxTeilkapitelAASS}}

{\sffamily\noindent Hier finden sich einzelne kurze
Abhandlungen zu Themen, welche im jeweiligen Kompetenzlevel
(noch) keinem Kapitel zugeordnet sind.}

\clearpage % TEMP temp clearpage

\section{Vorhofflimmern}
% \TextSlide
% 	{<Titel>}%1
% 	{<Text>}%2
% 	{<Slide>}%3
\TextSlide
{{\TxBeschreibung} und Ursachen}%1
{}%2
{}%3
Beim Vorhofflimmern kommt es zu
unkoordinierter Erregung der Vorhöfe durch 
\EE{Mikro-Reentrys} in den Vorhöfen bzw. an den 
Einmündung der \EE{Pulmonalvenen}.

VHF ist meist sekundär
bedingt: \EE{Kardiale Ursachen} umfassen u.\,a.
ischämische und entzündliche Herzerkrankungen,
Kardiomyopathien, Herzinsuffizienz und das Sick Sinus
Syndrome.
\E{Extrakardiale Ursachen} umfassen Systemerkrankungen wie
die arterielle Hypertonie, Lungenembolie, Hyperthyreose
oder toxische Wirkungen, z.\,B. durch Alkohol (\I{holiday
heart syndrome}) oder Medikamente (Thyroxin, Betamimetika
u.\,v.\,a.\,m.). VHF kann auch \E{primär} als \I{lone
atrial fibrillation} auftreten. 


% \TextSlide
% 	{<Titel>}%1
% 	{<Text>}%2
% 	{<Slide>}%3
\TextSlide
{Auftreten
  und
  resultierende
  Probleme}%1
{}%2
{}%3  
VHF kann selbstlimitierend \EE{paroxysmal} (<\,7 Tage),
anhaltend über mehr als 7 Tage \EE{persistierend} (wobei
eine Konversion in einen Sinusrhythmus möglich ist) oder
\EE{permanent} (ohne Möglichkeit der Konversion)
auftreten.
Es
ergeben sich folgende \EE{Probleme}:
\begin{itemize}
  \item \EE{Herzinsuffizienz}: Der Herzauswurf kann beim
  Gesunden um 15\PC, beim bereits Herzkranken um 40\PC
  vermindert sein.
  \item
  \EE{Thromben}: Durch den faktischen Stillstand der
  Vorhöfe können sich Thromben bilden, welche Embolien
  (Schlaganfälle, Mesenterialinfarkte, Lungenembolien, \ldots)
  verursachen können. Das Risiken für Cerbro-Vascular
  Incidents (CVI) sowie für Komplikationen bei
  prophylaktischer Antikoagulation können mit Scores
  abgeschätzt werden (s.\,u.).
\end{itemize}


% \TextSlide
% 	{<Titel>}%1
% 	{<Text>}%2
% 	{<Slide>}%3
\TextSlide
{Klinik}%1
{%
Die Palette der Symptome reicht von asymptomatisch bis zu
akuter kardiozirkulatorscher Instabilität und sind sehr von
den individuellen Auswirkungen des VHF abhängig (s.\,o.).

Patienten berichten mitunter über Palpitationen oder (prä-) synkopale Episoden, sowie unspezifische Beschwerden, wie Übelkeit und Schwindel,
Leistungsintoleranz, evtl. auch Dyspnoe, Stenokardien und
Thoraxschmerzen.}%2
{}%3



~~

\includegraphics[width=\linewidth]{images/Afib_ecg.jpg}


% \TextSlide
% 	{<Titel>}%1
% 	{<Text>}%2
% 	{<Slide>}%3
\TextSlide
{Diagnostik}%1
{%
\begin{itemize}
  \item Anamnese + physikalische Untersuchung
  \item Echokardiographie: Strukturell gesund?
  \item Labor (TSH!)
  \item Ggfs. Ergometrie, Koronarangiographie
  \item Holter-EKG
  \item Ggfs. TEE
\end{itemize}}%2
{%
\begin{itemize}
  \item Anamnese + physikalische Untersuchung
  \item Echokardiographie: Strukturell gesund?
  \item Labor (TSH!)
  \item Ggfs. Ergometrie, Koronarangiographie
  \item Holter-EKG
  \item Ggfs. TEE
\end{itemize}}%3


% \TextSlide
% 	{<Titel>}%1
% 	{<Text>}%2
% 	{<Slide>}%3
\TextSlide
{Antikoagulation ?}%1
{}%2
{}%3

\begin{tabulary}{\linewidth}{ccL}
\toprule\addlinespace
\multicolumn{3}{l}{\TabHeading{CHA\textsubscript{2}DS\textsubscript{2}-VASc} -- Thromboembolisches Risko}
\\\addlinespace\midrule\addlinespace
{\bfseries\color{AASSRed}C}&1&Congestive heart failure
\\
{\bfseries\color{AASSRed}H}&1&Hypertension
\\
{\bfseries\color{AASSRed}A}&{\bfseries\color{AASSRed}2}&Age > 75\,a
\\
{\bfseries\color{AASSRed}D}&1&Diabetes mellitus
\\
{\bfseries\color{AASSRed}S}&{\bfseries\color{AASSRed}2}&Stroke / TIA
\\
{\bfseries\color{AASSRed}V}&1&ascular disease
\\
{\bfseries\color{AASSRed}A}&1&ge \VonBis{65}{74}\,a
\\
{\bfseries\color{AASSRed}Sc}&1&Sex category: female
\\\addlinespace
\multicolumn{3}{l}{\small$\geq$\,2: OAK empf; =\,1: OAK od. ASS; =\,0: Evtl. ASS}
\\\addlinespace\toprule\addlinespace
\multicolumn{3}{l}{\TabHeading{HAS-BLED} -- Blutungsrisiko}
\\\addlinespace\midrule\addlinespace
{\bfseries\color{AASSRed}H} & 1 &  Hypertension 
   \\
{\bfseries\color{AASSRed}A}& 1 &  Abnormal renal \ldots 
\\
   & 1 &  {\ldots} liver function
\\
{\bfseries\color{AASSRed}S}& 1 &  Stroke
\\
{\bfseries\color{AASSRed}B}& 1 &  Bleeding 
\\
{\bfseries\color{AASSRed}L}& 1 &  Labile INRs 
\\
{\bfseries\color{AASSRed}E}& 1 &  Elderly ($\geq$ 65\,a) 
\\
{\bfseries\color{AASSRed}D} & 1 & Drugs or \ldots
\\
  & 1 &  {\ldots} Alcohol  
\\\addlinespace
\multicolumn{3}{l}{\small$\geq$\,3: High risk}
\\\addlinespace\bottomrule\addlinespace
\end{tabulary}

% CHA\textsubscript{2}DS\textsubscript{2}-VASc-Punkte

\begin{tabular}{cc}
\addlinespace\toprule\addlinespace
\TabHeading{CHA\textsubscript{2}DS\textsubscript{2}-VASc}
& 
\TabHeading{CVI-Risiko p.\,a. [\PC]}
\\\addlinespace\midrule\addlinespace
0	&	0\\
1	& 	1,3\\
2	&	2,2\\
\ldots	&	\ldots\\
6	&	9,8
\\
\ldots	&	\ldots\\
9 & 15,2
\\\addlinespace\bottomrule
\end{tabular}



% \TextSlide
% 	{<Titel>}%1
% 	{<Text>}%2
% 	{<Slide>}%3
\TextSlide
{De
  novo:
  Kardioversion
  (CV)}%1
{%
\EE{Innerhalb von 48\Hour*} kann eine primär
\EE{medikamentöse CV} erfolgen. \EE{Nach 48\Hour*} muss erst
für mind. 3 Wochen eine suffiziente orale Antikoagulation
(OAK) erfolgen. In beiden Fällen kann bei misslingen der
medikamentösen CV eine \EE{elektrische CV} durchgeführt
werden. OAK mind. 4 Wochen weiter!

\subparagraph{Medikamentöse CV} je nach Zustand
des Herzens:
\begin{itemize}
  \item Herzkrank: Sedacoron
  \item Herzgesund: Propafenon (Rytmonorma{\texttrademark}),
  Ibutilide (Covert{\texttrademark}), Flecainid
  (Aristocor{\texttrademark})
\end{itemize}
\bisgskip 

\subparagraph{Elektrische Kardioversion} wird in Kurznarkose
und R-Zacken-\Ca{synchronisiert!} durchgeführt.


\subparagraph{Rezidivprophylaxe} nach dem ersten Rezidiv:
\begin{itemize}
  \item Herzkrank (strukturell): Sedacoron
  \item KHK: Sotalol (Sotacor{\texttrademark},
  Beta-Blocker-Wirkung)
  \item Herzgesund: Propafenon (Rytmonorma{\texttrademark}),
  Flecainid
  (Aristocor{\texttrademark})
\end{itemize}}%2
{}%3


% \TextSlide
% 	{<Titel>}%1
% 	{<Text>}%2
% 	{<Slide>}%3
\TextSlide
{Intervention}%1
{%
Weiters besteht bei
\E{symptomatischen} VHF die Möglichkeit der
Pulmonalvenenisolation (PVI): 
Dabei werden durch Ablation die Einmündungen der Pulmonalvenen elektrisch isoliert.
Erfolg nach 1 Jahr: \VonBis{30}{50}\PC}%2
{}%3



% \begin{compactenum}
%   \item < 12 Monate
%   \item keine kardiale Grunderkrankung
%   \item Linker VH < 50\Mm* Durchmesser
%   \item Kausale Ursachen beseitigt
%   \item Kein SS
%   \item Mitralfehler $\leq$ II
%   \item Kardiale Grunderkrankung: Monitorkontrolle
%   \item Thromben ausgeschlossen (< 48\Hour*, TEE,
%   OAK \VonBis{4}{5} Wochen im Zielbereich)
% \end{compactenum}

% 
% 
% \paragraph{Flecainid (Aristocor{\texttrademark}):}
% \VonBis{200}{300}\Mg*; \E{"`Pill in the pocket"'}; KI:
% Strukturelle Herzkrankheit --
% \Ca{Cave: Maligne Arrhythmien bei Klasse-I-Antiarrhythmika!}
% 
% 
% 
% \paragraph{Amiodaron}
% \textbf{{Rp.}} Loading dose 200\,Mg* 1-1-1 p.\,o. für 10
% Tage, dann 1-0-0 
% 
%   

% \TextSlide
% 	{<Titel>}%1
% 	{<Text>}%2
% 	{<Slide>}%3
\TextSlide
{}%1
{}%2
{}%3

Persistierend:
  Frequenzkontrolle
  

\TabHeading{Verapamil} oder Diltiazem kann bei Patienten \emph{ohne}
Herzinsuffizienz eingesetzt werden. \Ca{Cave: Bei Kombination mit Beta-Blocker AV-Block!}

\TabHeading{Digitalis} kann als Ergänzung bei
Herzinsuffizienz (in Kombination mit Beta-Blockern) eingesetzt werden; wirkt
negativ chronotrop und positiv inotrop.

\TabHeading{Beta-Blocker} eignen sich besonders bei
Herzinsuffizienz und bei Hyperthyreose



% \TextSlide
% 	{<Titel>}%1
% 	{<Text>}%2
% 	{<Slide>}%3
\TextSlide
{}%1
{}%2
{}%3



% \paragraph{Digoxin} alternativ bei systolischer
% Herzinsuffizienz 
% \EE{Rp.} 0,25\Mg* \Abk{iv} in KI, bei
% Bedarf 2-stündlich bis max.
% 1,25\Mg*\,/\,d

% \TextSlide
% 	{<Titel>}%1
% 	{<Text>}%2
% 	{<Slide>}%3
\TextSlide
{Bradyarrhythmia absoluta}%1
{VVI(R)-Schrittmacher mit frequenzadaptierter Stimulation

}%2
{}%3

% \TextSlide
% 	{<Titel>}%1
% 	{<Text>}%2
% 	{<Slide>}%3
\TextSlide
{Antikoagulation!}%1
{\TabHeading{Coumarine -- First line}
Erfolgskontrolle über INR (oder Quick)  möglich, Antidot
verfügbar.

Risiko einer
intrakraniellen Blutung 0,3\PC p.\,a.


\TabHeading{Alternative OAK} z.\,B. Pradaxa, Xarelto: Kontroversiell diskutiert. Keine
INR-Kontrolle notwendig, jedoch auch keine Erfolgskontrolle
möglich. Derzeit Verschreibung nur mittels IND.

\TabHeading{ASS} Bei absoluter Kontraindikation zur OAK oder
CHADS\textsubscript{2}-Score\,=\,1 und kein Benefit von OAK. Evtl. auch bei CHADS\textsubscript{2}-Score\,=\,0
}%2
{}%3





\ChapterEnd